\section{Model and results}
\label{sec:model}

Reinforcement learning with verifiable rewards (RLVR) raises a base model's
pass@1, yet at large sampling budgets $k$ the base model catches up and
overtakes it. This reversal has been read as evidence that RLVR only sharpens
existing ability rather than adding reasoning (Yue et al., 2504.13837, ``Does
RL Really Incentivize Reasoning''). We give it an exact reading: the reversal is
a \emph{price}, and we compute its exchange rate. Modeling RLVR as an
exponential tilt over strategies shared across prompts, we prove that once
strategies specialize, every update that raises pass@1 must lower the success of
some prompt, with a theorem of the alternative supplying the exact condition (Theorem
\ref{thm:nfl}); that when the loss reaches the hardest prompts it forces the
pass@$k$ crossover, whose number is bounded by the sign changes of the shift in
per-prompt success (Theorem \ref{thm:price}); and that a constrained update
traces the optimal pass@1 versus pass@$k$ frontier (Theorem \ref{thm:frontier}).
There is no free lunch, only a price. We state the model and all results here;
proofs are in the appendix, and every numerical constant is produced and checked
by the accompanying package.

\paragraph{Setup.}
A base policy is refined by reinforcement learning with verifiable rewards
(RLVR). We model the object that RLVR moves as a distribution over
\emph{solution strategies} that are shared across prompts. There are $M$
strategies and a prompt distribution $\mathcal{D}$ on a finite set of $N$
prompts. For a prompt $x$ and a strategy $j\in[M]$, let
\begin{equation}
q_{xj}\in[0,1]
\quad\text{be the probability that strategy $j$ solves prompt $x$,}
\qquad Q=(q_{xj})\in[0,1]^{N\times M}.
\end{equation}
A policy is a point $w\in\simplex$ on the strategy simplex. Drawing one strategy
$j\sim w$ and solving $x$ with probability $q_{xj}$, the per-prompt success
probability and the mean (pass@1) success are
\begin{equation}
p_x(w)=\inner{w}{q_x}=\sum_{j}w_j q_{xj},
\qquad
\pass@1(w)=\E_{x\sim\mathcal{D}}\,p_x(w)=\inner{w}{f},
\quad
f_j:=\E_{x\sim\mathcal{D}}\,q_{xj}.
\end{equation}
Here $q_x=(q_{x1},\dots,q_{xM})$ is the $x$-th row of $Q$ and $f=\E_x q_x$ is the
vector of \emph{marginal fitnesses}. With $k$ independent samples from the
policy, the pass@$k$ of prompt $x$ and its average are
\begin{equation}
\pass@k(x;w)=1-(1-p_x(w))^k,
\qquad
\PassK_k(w)=\E_{x}\big[1-(1-p_x(w))^k\big].
\end{equation}
We write $s_x=-\log(1-p_x)\in[0,\infty)$; hard prompts have small $p_x$, hence
small $s_x$, and the large-$k$ behavior of $\PassK_k$ is governed by the small-$s$
(hardest) prompts through $(1-p)^k=e^{-ks}$.

\paragraph{The RLVR update.}
GRPO and its relatives subtract a group baseline and rescale by the group
standard deviation, then push probability toward above-baseline strategies. In
the population limit, with a KL anchor to the base, the idealized update is
mirror descent / natural policy gradient on the strategy logits. Theorem
\ref{thm:tilt} identifies the resulting policy as an exponential tilt of the
base. Throughout, $\wz\in\mathrm{int}\,\simplex$ denotes the base policy and
$\wt$ the policy after internal training time $\tau\ge 0$.

\paragraph{Comparison objects.}
For a base policy $\wz$ and a trained policy $\wt$, write
$F_0,F_\tau$ for the cumulative distribution functions of the per-prompt
success values $\{p_x(\wz)\}$ and $\{p_x(\wt)\}$ under $\mathcal{D}$, and
\begin{equation}
G(p)=F_\tau(p)-F_0(p),
\qquad
\Delta_k=\PassK_k(\wt)-\PassK_k(\wz).
\end{equation}
A \emph{pass@$k$ crossover} is a sign change of $k\mapsto\Delta_k$. We write
$\Svar(\cdot)$ for the number of sign changes of a sequence or function.

\subsection{Assumptions}

\begin{assumption}[Exclusive specialization]
\label{ass:spec}
Each strategy is the sole solver of a set of prompts of positive mass: for each
$j\in[M]$ there is a set $X_j^{\mathrm{excl}}$ with
$\mathcal{D}(X_j^{\mathrm{excl}})>0$ such that $q_{xj}>0$ and $q_{xj'}=0$ for all
$j'\ne j$ on $X_j^{\mathrm{excl}}$.
\end{assumption}

Exclusivity, not mere being-best, is what makes probability a scarce resource:
moving mass off $j$ strictly lowers the prompts only $j$ can solve. Being the
\emph{unique best} is \emph{not} enough, because a near-generalist can absorb the
mass without hurting any prompt; we record this as Proposition
\ref{prop:uniquebest}. A realistic version without exact zeros, used in the
experiments, is dominant specialization: for each $j$ a positive-mass set with
$q_{xj}\ge d$ and $q_{xj'}\le\varepsilon$ for $j'\ne j$, with
$(M-1)\varepsilon<d/2$; it is sufficient (Theorem \ref{thm:nfl}(b$'$)).

\begin{assumption}[Fine coverage]
\label{ass:fine}
The prompts of minimal base success $x^\star\in\arg\min_x p_x(\wz)$ are covered
only by strategies of below-base-average fitness: $q_{x^\star j}>0\Rightarrow
f_j<R_0$, where $R_0=\pass@1(\wz)=\inner{\wz}{f}$ is the base pass@1.
\end{assumption}

Because the tilt multiplies $w_{0,j}$ by $e^{\tau(f_j-\Lambda(\tau)/\tau)}$ with
$\Lambda(\tau)=\log\E_{\wz}e^{\tau f}$, a strategy loses mass iff
$f_j<\Lambda(\tau)/\tau$, and Jensen gives $\Lambda(\tau)/\tau\ge R_0$ for all
$\tau>0$. Hence under Assumption \ref{ass:fine} every strategy covering $x^\star$
loses mass, so $p_{x^\star}(\wt)<p_{x^\star}(\wz)=\min_x p_x(\wz)$. (The looser
threshold $f_j<R(\tau)$, the \emph{tilted} mean, is wrong: a strategy with
$R_0<f_j<\Lambda(\tau)/\tau$ gains mass though $f_j<R(\tau)$; see the appendix.)
Assumption \ref{ass:spec} forces a per-prompt loss (Theorem \ref{thm:nfl});
Assumption \ref{ass:fine} routes it to the hardest prompt and thereby forces an
aggregate crossover (Theorem \ref{thm:price}). The two must not be conflated:
the per-prompt loss is unconditional under \ref{ass:spec}, the crossover is
conditional on \ref{ass:fine}, and \ref{thm:price}(e) is the regime where
\ref{ass:fine} fails and no crossover forms.

\subsection{Statements}

\begin{theorem}[Tilt identity]
\label{thm:tilt}
Under mirror descent / natural policy gradient on the strategy logits with a KL
anchor of strength $\beta$ and binary verifiable reward, the policy at internal
time $\tau$ is the exponential tilt
\begin{equation}
\wt \;\propto\; \wz\,\exp(\tau f),
\qquad
w_{\tau,j}=\frac{w_{0,j}\,e^{\tau f_j}}{\sum_{\ell}w_{0,\ell}\,e^{\tau f_\ell}} .
\end{equation}
This is exact for mirror descent / natural policy gradient.
\end{theorem}

\begin{proposition}[GRPO stays near the tilt]
\label{prop:grpo}
For GRPO with the group non-degeneracy factor $1-p^G-(1-p)^G$, the policy stays
within total variation $O(\delta)$ of the best-fit tilt at matched concentration,
where $\delta$ is the variation of the group standard-deviation rescaling over the
trajectory. We do not optimize the constant; the package measures this total
variation below $3\times10^{-4}$ over the tested range, and records the
finite-group transient non-monotonicity separately.
\end{proposition}

\begin{theorem}[No free lunch]
\label{thm:nfl}
A \emph{free lunch} at an interior $w$ is a tangent $\mathrm{d}w$
($\mathbf 1^\top\mathrm{d}w=0$) with $Q\,\mathrm{d}w\ge0$ and $Q\,\mathrm{d}w\ne0$;
since $\inner f{\mathrm{d}w}=\tfrac1N\mathbf 1^\top Q\,\mathrm{d}w$, this is a
pass@1 gain with no prompt regressing. Whether one exists does not depend on $w$.
\begin{enumerate}
\item[(a)] \emph{(Exact condition, Stiemke.)} No free lunch exists if and only if
there are strictly positive prompt prices $y\in\R^{N}_{>0}$ and $\mu\in\R$ with
\begin{equation}
\label{eq:farkas}
\sum_x y_x\,q_{xj}=\mu\qquad\text{for all } j\in[M],
\end{equation}
a strictly positive reweighting of prompts under which all strategies have equal
weighted fitness. (Equivalently, since $f=\tfrac1N\sum_x q_x$, there exist
$y'\ge0,\mu'$ with $f_j+\sum_x y'_x q_{xj}=\mu'$.)
\item[(b)] \emph{(Exclusive specialization suffices.)} Under Assumption
\ref{ass:spec} no free lunch exists: any $\mathrm{d}w$ with $\mathrm{d}w_j<0$
strictly lowers $X_j^{\mathrm{excl}}$, so every pass@1 gain lowers some prompt.
\item[(b$'$)] \emph{(Dominant specialization suffices.)} If for each $j$ there is a
positive-mass set with $q_{xj}\ge d$ and $q_{xj'}\le\varepsilon$ ($j'\ne j$) and
$(M-1)\varepsilon<d/2$, then \eqref{eq:farkas} is feasible with $y>0$ and no free
lunch exists.
\item[(c)] \emph{(Necessity.)} If some strategy dominates all others
($q_{xj^\star}\ge q_{xj}$ for all $x,j$), concentrating on it raises every $p_x$:
a free lunch.
\end{enumerate}
\end{theorem}

\begin{proposition}[Unique-best is insufficient]
\label{prop:uniquebest}
Being the unique best scorer on a positive-mass set does not prevent a free lunch.
For $Q=\begin{psmallmatrix}1&0.9&0.96\\0.9&1&0.96\\0.5&0.5&1\end{psmallmatrix}$
each strategy is the unique best on its prompt, yet
$\mathrm{d}w=(-0.6,-0.4,1)$ gives $Q\,\mathrm{d}w=(0,0.02,0.5)\ge0$ and
$\inner f{\mathrm{d}w}=+0.173$: a free lunch, and \eqref{eq:farkas} is infeasible.
\end{proposition}

\begin{theorem}[The pass@$k$ price]
\label{thm:price}
Let $\wz,\wt$ be a base and a trained policy with success-CDF difference $G$.
\begin{enumerate}
\item[(a)] \emph{(Crossing bound.)} The number of pass@$k$ crossings is bounded
by the sign variation of $G$:
\begin{equation}
\Svar\big(k\mapsto\Delta_k\big)\;\le\;\Svar(G),
\end{equation}
because $\Delta_k=-k\int_0^1 (1-p)^{k-1}G(p)\,\mathrm{d}p$ and the kernel
$k(1-p)^{k-1}$ is totally positive in $(k,p)$ (Karlin variation diminishing).
\item[(b)] \emph{(Tightness.)} For every $m\ge 0$ there is a configuration with
exactly $m$ crossings; the bound is attained.
\item[(c)] \emph{(Forced crossover.)} Under Assumptions \ref{ass:spec} and
\ref{ass:fine}, any pass@1-increasing tilt has $p_{x^\star}(\wt)<\min_x p_x(\wz)$,
hence $\Delta_k<0$ for all sufficiently large $k$: a crossover is forced, and its
count obeys (a).
\item[(d)] \emph{(Uniqueness; recovers prior work.)} If the trained policy
preserves the difficulty order of prompts, $p_x(\wt)=\phi(p_x(\wz))$ for an
increasing $\phi$, then $\Svar(G)\le 1$ and there is at most one crossing. This
recovers the single-crossover regime of Yang and Chen (2609.22547, Thm.~3) as a
corollary.
\item[(e)] \emph{(Counterexample.)} Without Assumption \ref{ass:fine}, the loss
need not reach the hard tail and no crossover occurs, even though a per-prompt
loss does; this regime bounds the scope of (c) and makes the result
falsifiable.
\end{enumerate}
\end{theorem}

\begin{theorem}[Price frontier]
\label{thm:frontier}
For $\gamma\in[0,1]$, let $w(\gamma)$ maximize $\pass@1$ subject to the
per-prompt floor $p_x(w)\ge\gamma\,p_x(\wz)$ for all $x$. The map
$\gamma\mapsto\big(\pass@1\text{ gain},\,k^\star\big)$ is monotone: $\gamma=0$
recovers the unconstrained gain with the earliest crossing, and $\gamma=1$
enforces first-order stochastic dominance $F_\tau\le F_0$, hence $\Delta_k\ge0$
for all $k$ and zero pass@1 gain under Assumption \ref{ass:spec}. The frontier
is the exact exchange rate between pass@1 and pass@$k$; it is a linear program
in $Q$ and has a closed form for $M=2$.
\end{theorem}

\begin{theorem}[Tilt law]
\label{thm:tiltlaw}
Let the base be uniform and let the fitness density near its maximum $b$ behave
as $c\,g^{\kappa-1}$ in the gap $g=b-f$. In the mesoscopic regime the
entropy-performance relation is
\begin{equation}
b-R \;=\; a_\kappa\,\exp(H/\kappa)\,(1+o(1)),
\end{equation}
where $R=\inner{\wt}{f}$ and $H=H(\wt)$. The exponent $\kappa$ is identifiable
from the slope of $\log(b-R)$ against $H$; $\kappa=1$ recovers
$R=-a\,e^{H}+b$ of Cui et al. (2505.22617). This result is stated for context;
the exact entropy-mass curve and the non-exponential terminal law are due to
concurrent work (Winner-Take-All, S4R2zYtFhs).
\end{theorem}

\begin{lemma}[Entropy bump]
\label{lem:bump}
Along the tilt, $\tfrac{\mathrm{d}}{\mathrm{d}\tau}H(\wt)
=-\mathrm{Cov}_{\wt}(f,\log \wz)-\tau\,\mathrm{Var}_{\wt}(f)$.
Entropy rises before it collapses if and only if
$\mathrm{Cov}_{\wz}(f,\log \wz)<0$, that is, the high-fitness strategies are
initially underweighted.
\end{lemma}

\begin{lemma}[Sharing is necessary]
\label{lem:sharing}
If each prompt is served by its own independent (tabular) policy and the update
follows the full gradient flow, then $p_x$ is nondecreasing for every $x$, so
$\Delta_k\ge0$ for all $k$ and no inversion occurs. A pass@$k$ inversion
therefore requires strategies shared across prompts. Under finite-group
stochastic GRPO the monotonicity can fail transiently (2602.06717).
\end{lemma}

\paragraph{What is new, and why it matters.}
Theorem \ref{thm:nfl} is the headline and is unconditional given Assumption
\ref{ass:spec}: concentration buys pass@1 only by spending some prompt's
success, and the pricing identity \eqref{eq:farkas} values the trade exactly.
This answers, with a mechanism, why RLVR cannot widen the reasoning frontier
for free and says what it costs when it improves pass@1. Theorem
\ref{thm:price}(a) bounds the number of pass@$k$ crossings for an arbitrary
shift in per-prompt success; this is the general statement that concurrent work
leaves open for kernels and multiple crossings (Yang and Chen, 2609.22547,
Rem.~1), and it recovers their single-crossover theorem as the order-preserving
corollary \ref{thm:price}(d). Theorem \ref{thm:price}(c) then turns the
per-prompt price into the observed aggregate phenomenon, the crossover, under
the named Fine Coverage condition, with the counterexample regime
\ref{thm:price}(e) bounding its scope and making it falsifiable. Theorem
\ref{thm:frontier} is the exact exchange rate and the only place a
``correction'' appears: not a free fix, since none exists, but the frontier that
attains the minimal price. Theorem \ref{thm:tiltlaw} and the lemmas tie the same
tilt identity to the entropy side studied in concurrent work. The practical
read-off is a diagnostic that takes only per-prompt base and RL success rates,
available in any evaluation log, and returns the crossing bound, the
order-preserving test, and the estimated price. There is no free lunch, only a
price.